\icode{N24J-ME-MU-L3-261019-01}
\nceeID[6]{2024吉林}{10}一足够长木板置于水平地面上，二者间的动摩擦因数为\(\mu\)。\(t = 0\)时，木板在水平恒力作用下，由静止开始向右运动。某时刻，一小物块以与木板等大、反向的速度从右端滑上木板。已知\(t = 0\)到\(t = 4t_0\)的时间内，木板速度\(v\)随时间\(t\)变化的图像如图所示，其中\(g\)为重力加速度大小。\(t = 4t_0\)时刻，小物块和木板的速度相同。下列说	法正确的是\paren[]

\begin{tikzpicture}[>=stealth,font=\scriptsize]
	\draw[->,blue] (0,0)node[below left]{\(O\)}--(4,0)node[below]{\(t\)};
	\draw[->,blue] (0,0)--(0,3)node[right]{\(v\)};
	\foreach \x/\n in{.8/1,1.6/2,2.4/3,3.2/4}
		\draw[blue] (\x,.05)--(\x,0)node[below,blue]{\(\n t_0\)};
	\draw[dashed,cyan] (0,.7)node[left,blue]{\(\frac{1}{2}\mu gt_0\)} --(3.2,.7)--(3.2,0) (.8,.7)--(.8,0) (2.4,2.1)--(2.4,0);
	\draw[magenta] (0,0)--(2.4,2.1)--(3.2,.7);
\end{tikzpicture}
\begin{choices}
	\item 小物块在\(t = 3t_0\)时刻滑上木板
	\item 小物块和木板间的动摩擦因数为\(2\mu\)
	\item 小物块与木板的质量比为\(3:4\)
	\item\(t = 4t_0\)之后小物块和木板一起做匀速运动
\end{choices}
\begin{answer}
	
\end{answer}